\chapter{Instruction Selection}\label{ch:instruction-selection}

Instruction selection turns machine-independent IR into machine instructions.
The IR says what to compute without saying how the target spells it.
The selector closes that gap by covering each IR fragment with target patterns that compute the same value.
Good coverage produces short, fast code.
Bad coverage produces correct but clumsy code with redundant moves, missed addressing modes, and unused compound instructions.

\section{Tiles, Patterns, and Costs}\label{sec:ch25-tiles}

A selector views the program as trees or graphs and the machine as a catalog of patterns.
\term{Instruction selection} maps IR operations onto target instructions that implement them.
\term{Pattern} means a template IR fragment with a result, a cost, and an emission recipe.
\term{Tile} means one placed pattern instance covering part of the program.
\term{Cover} means a set of tiles that implements every operation exactly once without overlap or gaps.
The task looks like jigsaw: cover the IR with the cheapest fitting tiles, then emit their recipes.

Consider address computation on a target with scaled indexed addressing.
The IR fragment \code{t = add x, mul k, 8} followed by \code{v = load [t]} can become one load with a scaled index when \code{k} is the loop counter.
Selected separately, the same fragment becomes a multiply, an add, and a plain load.
The difference is entirely in the selector window: a peephole of one operation misses the compound form, while a tree window of three operations finds it.
That is why selectors work on expression trees or directed acyclic graphs rather than single instructions.
The \code{load}, \code{mul}, and \code{add} in the example are \code{IR opcodes} that one load-with-scale tile absorbs.

Costs guide the choice when several covers are legal.
A simple model assigns each pattern a static size or latency number.
A better model adds context such as whether an operand is a constant, whether a value is already in the right register class, and whether the addressing mode folds a use for free.
For a fixed finite catalog with additive local costs, tree dynamic programming finds a minimum-cost cover in linear time in tree size.
This optimum is relative to the catalog and cost model; register pressure, scheduling, and shared subexpressions introduce interactions outside that model.
The practical consequence is modest: tree selectors with a sharing rule handle nearly all integer and address code well, and the remaining graph cases get careful hand patterns.
Table~\ref{tab:ch25-pattern-examples} lists typical patterns and what they absorb.
Figure~\ref{fig:ch25-tiling} shows two competing covers of one expression tree.

\begin{figure}[htbp]
\centering
\begin{tikzpicture}[node distance=0.9cm and 1.2cm]
\node[basic block] (add) {add};
\node[basic block, below left=of add] (mul) {mul k, 8};
\node[basic block, below right=of add] (x) {x};
\draw[cfg edge] (add) -- (mul);
\draw[cfg edge] (add) -- (x);
\end{tikzpicture}
\caption{The multiply and add form an address subtree. When consumed by a load, the whole fragment may match a scaled-index load tile.}
\label{fig:ch25-tiling}
\end{figure}

\begin{table}[htbp]
\caption{Common selection patterns and the IR they absorb.}
\label{tab:ch25-pattern-examples}
\begin{tabular}{lll}
\toprule
Pattern & IR fragment absorbed & Benefit \\
\midrule
Load with displacement & \code{add} plus \code{load} & Fewer instructions \\
Scaled indexed load & \code{mul}, \code{add}, \code{load} & No separate multiply \\
Multiply-accumulate & \code{mul} plus \code{add} & One fused operation \\
Compare and branch & \code{cmp} plus \code{br} & Uses flag-setting form \\
Move immediate & Constant plus \code{mov} & Shorter encoding \\
\bottomrule
\end{tabular}
\end{table}

The target description supplies the pattern catalog.
Each rule names the IR shape, the constraints on operands, the result location, and the assembly or machine-IR template.
Predicates guard rules that apply only to certain constant ranges, alignments, or register classes.
A RISC target contributes a few dozen crisp rules.
A vector or DSP target contributes hundreds with overlapping shapes.
Keeping the catalog ordered and tested matters more than making it large.
An untested rule with a wrong predicate silently miscompiles every program that triggers it.
\begin{ccwarning}[Every pattern is a correctness claim]
A pattern asserts semantic equivalence between its IR fragment and its instruction template for all allowed operands.
Test constant boundaries, signedness, overflow flags, and NaN behavior explicitly.
A missed edge case in one rule corrupts every function the selector tiles with it.
\end{ccwarning}

\section{Maximal Munch and Tree DP}\label{sec:ch25-techniques}

Two classical techniques cover most practice: greedy \term{maximal munch} and optimal tree dynamic programming.
Maximal munch walks the tree from the root and always takes the largest pattern that matches at the current node, then recurses on the leftovers.
It is fast, easy to implement, and good enough for regular RISC targets where the largest tile is usually the best tile.
Its weakness is myopia: it can grab a large tile at the root that forces awkward tiles below, while two smaller tiles would have been cheaper overall.
Tree DP fixes that by computing the cheapest cover for every subtree bottom-up and then emitting top-down.
Each node records the best cost per result location, called its \term{state}, and the rule that achieved it.
Emission follows the recorded rules without further search.

DAG structure complicates both techniques because a shared value has two parents.
A simple strategy cuts at shared values, materializes each once, and selects the resulting trees with temporary leaves.
Duplicating a shared computation is a separate legality and profitability decision, especially for loads, traps, and effects.
More ambitious selectors handle small DAGs directly or use multi-output patterns, such as an instruction producing both quotient and remainder.
Control flow stays outside the tile: selectors cover single-entry single-exit regions such as basic blocks or extended blocks, and branches become explicit tile boundaries.
Machine idioms such as zeroing a register with \code{xor r, r} or testing with a flag-setting compare belong in the catalog as special low-cost rules.

Algorithm~\ref{alg:ch25-tree-dp} gives the tree DP skeleton used in many textbook and production selectors.
It assumes each IR region has been converted to expression trees with explicit temporaries at region boundaries.
It labels every node with its cheapest rule per result kind, then emits code greedily from the root following those labels.
The labeling pass visits each node once and tries every rule whose shape matches, so the cost is $O(N \cdot R \cdot P)$, where $N$ is the node count, $R$ is the number of result kinds, and $P$ is the pattern count checked per node.
With indexed rule tables this behaves like linear time in practice.
Emission is linear in the tile count.

\begin{algorithm}[htbp]
\caption{Tree-pattern instruction selection by dynamic programming}\label{alg:ch25-tree-dp}
\begin{pseudocode}
\Require expression tree $T$, finite tree patterns $P$ with additive costs and proper-subtree frontier operands
\Ensure cheapest cover for the requested root result kind, or failure if none exists
\ForAll{node $n$ in postorder}
\State set $Cost(n, k) \gets \infty$ for each result kind $k$
\ForAll{rule $r$ in $P$ matching $n$}
\State check operand predicates and constant ranges
\State $c \gets Cost(r) + \sum_{f\in frontier(r,n)} Cost(f, kind(r,f))$
\If{$c < Cost(n, kind(r))$}
\State record $Cost(n, kind(r)) \gets c$ with rule $r$
\EndIf
\EndFor
\EndFor
\State reject if the required root cost is infinite
\State follow recorded rules, emitting frontier operands before their consuming instruction
\end{pseudocode}
\end{algorithm}

The frontier consists of unmatched operand subtrees at a tile's boundary, not necessarily the root's immediate children.
Charging covered internal nodes again would double-count the operations absorbed by a compound instruction.
The recurrence excludes same-node chain rules; catalogs with result-kind conversions need an additional closure computation.
The same skeleton supports single-node simplification patterns when their semantics and costs satisfy these conditions.
Constant folding at selection time, identity removal, and strength reduction by pattern all fall out of the cost comparison without a separate pass.
The danger is rule interaction: two simplification rules can ping-pong the same node if their costs tie, so implementations order rules deterministically and re-label to a fixed point only within one node.
Retargetability comes from generating the matcher from the catalog rather than hand-writing nested tests.
Table-driven matchers, tree-grammar generators, and declarative TableGen-style patterns all produce the same bottom-up labeling behavior with different engineering tradeoffs.
% TODO(cite): need a primary reference for tree-grammar code-generator generators.

\section{Worked Example}\label{sec:ch25-worked}

We tile a small address expression by hand on a target with three relevant patterns: plain load, load with 12-bit displacement, and scaled indexed load.
The program computes \code{a[i+1]} for 8-byte elements with base \code{a} and index \code{i}.
\begin{irlisting}
; IR trees
t1 = mul i, 8
t2 = add t1, 8
t3 = add a, t2
v = load t3
\end{irlisting}
Assign costs: plain operations cost 1 each, displacement folding saves 1, and the scaled indexed form costs 2 total.
Bottom-up labeling starts at \code{mul i, 8}: the scaled-index fragment matches with cost carried upward, while plain multiply costs 1.
At \code{add t1, 8}, the displacement rule matches because \code{8} fits in 12 bits, giving a folded cost below the two-instruction alternative.
At the root, the scaled indexed load covers the multiply, both adds, and the load in one tile of cost 2 plus a materialized base.
Count precisely: the naive cover uses \code{mul} (1), \code{add} (1), \code{add} (1), and \code{load} (1) for a total of 4.
The displacement-only cover uses a multiply, a base-plus-index add, and a displacement load, each of cost $1$, for a total of $3$; the full scaled form costs $2$.
The selector picks the cheapest, namely the scaled indexed load with the small constant folded into the displacement field.

Now verify legality of the winning tile step by step.
The multiply by 8 matches the scale field exactly, so no shift residue remains.
The constant 8 fits the displacement range, so no extra addend register is needed.
The base \code{a} is loop-invariant and already in a register, so the addressing mode can name it directly.
The loaded width is 8 bytes, which the mode supports.
Every predicate holds, so emission is one instruction plus the loop-invariant base setup.
If any predicate had failed, for example a 9000-byte displacement on a 12-bit field, the selector would fall back to the next-cheapest cover automatically.
That fallback chain is the point of costed DP over greedy matching.

\section{Implementation Notes}\label{sec:ch25-implementation}

Production selectors blend the textbook algorithm with target-specific lowering passes.
LLVM runs DAG combining before and during selection to canonicalize shapes toward the patterns the catalog favors, then uses a generated matcher over selection DAGs with hand-written C++ for complex cases.
GCC expands to a register-transfer level first and then matches machine-description patterns with predicates and constraints.
Both pipelines canonicalize commutative operations, fold constants early, and legalize types the target cannot handle directly.
Legalization matters: a 64-bit multiply on a 32-bit target becomes a library call or a multi-instruction sequence before tiling, not during it.

Debugging a selector means inspecting the DAG, the chosen tiles, and the predicate trace for the rejected tile.
Dump the selection DAG with costs per node when a suspicious instruction appears.
Check whether the missing tile failed on shape, predicate, type, or register class.
Common bugs include an off-by-one constant range, a signedness mismatch on extension patterns, and a missing predicate for an immediate that the assembler cannot encode.
Fuzz the catalog with random expression trees and differential-test against an unoptimized selector to catch miscompilations early.

Complexity stays manageable because regions are small.
Tree DP is linear in the region size for a fixed catalog, and DAG combining is bounded by worklist updates over local neighborhoods.
The engineering cost sits in the catalog: hundreds of patterns with predicates, costs, and tests.
Keep patterns declarative, review new rules for overlap with old ones, and record why each cost was chosen.
Good structural background on selector organization and machine-level IR appears in~\cite{appel-1998-modern} and see~\cite{morgan-1998-optimizing}.
Validation strategies for keeping lowering and selection faithful appear in~\cite{safonov-2010-trustworthy}.
The integer and type rules that predicates must respect are defined in~\cite{wg14-2011-c11}.
Chapter~\ref{ch:rewrite-optimization} places target patterns inside the wider rewrite framework, separates semantic validity from target cost, and tests a small multiply-add cover against explicit alternatives.
Chapter~\ref{ch:compiler-graph-transformation} explains why shared DAG rewrites also need preserved interfaces, use maintenance, and analysis invalidation rather than tree matching alone.

\section{Exercises}\label{sec:ch25-exercises}

\begin{enumerate}
\item Tile \code{t = add x, 16} followed by \code{v = load t} with costs 1 per plain operation and 1 for a displacement load. \label{ex:ch25-displacement} Show both covers, compute their total costs, and state the winning tile.
\item A shared value \code{s = mul p, 4} feeds two loads. \label{ex:ch25-sharing} Explain why tree duplication miscompiles or duplicates work here, and show the single-computation cover with explicit reuse.
\item A target adds a fused multiply-add with cost 2 versus cost 3 for separate multiply and add. \label{ex:ch25-fma-cost} Give an expression where maximal munch misses the fused form and tree DP finds it, with costs shown.
\end{enumerate}

\section{Further Reading}\label{sec:ch25-reading}

Selector structure, frame lowering, and machine-level code generation are introduced in~\cite{appel-1998-modern}.
Pattern design, machine descriptions, and optimization interplay are developed in~\cite{morgan-1998-optimizing}.
Techniques for validating aggressive low-level rewrites appear in~\cite{safonov-2010-trustworthy}.
Language-level semantics that selection must preserve are specified in~\cite{wg14-2011-c11}.
